\documentclass[11pt]{article}

% Change "review" to "final" to generate the final (sometimes called camera‑ready)
% version. Change to "preprint" to generate a non‑anonymous version with page numbers.
\usepackage[review]{acl}

% Standard package includes
\usepackage{times}
\usepackage{latexsym}

% For proper rendering and hyphenation of words containing Latin characters (including
% in bib files). See https://www.latex-project.org/help/documentation/encguide.pdf for
% other character sets
\usepackage[T1]{fontenc}

% This assumes your files are encoded as UTF8
\usepackage[utf8]{inputenc}

% This is not strictly necessary, and may be commented out, but it will improve the
% layout of the manuscript, and will typically save some space.
\usepackage{microtype}

% This is also not strictly necessary, and may be commented out. However, it will
% improve the aesthetics of text in the typewriter font.
\usepackage{inconsolata}

%Including images in your LaTeX document requires adding additional package(s)
\usepackage{graphicx}


\title{MultiIFEval: An Instruction-Following Resource for 305 Languages}

% No author names or acknowledgements in the review submission. Restore for camera-ready.
\author{Anonymous submission}

\begin{document}
\maketitle


\begin{abstract}

We present MultiIFEval, a publicly available translation and localisation of the
verifiable instruction-following benchmark IFEval \cite{zhou2023ifeval} into 305
language configurations. A structured generation pipeline adapts English-specific
constraints, then holds the resulting instruction identifiers fixed while translating
prompts and language-dependent constraint values. A Wikipedia article in the target
language provides contextual text. We describe the
released resource and its construction, and propose a focused zero-shot evaluation of
one or two small non-reasoning models on a subset of European languages using
EuroEval's instruction-following metric. This limited experiment is intended as a
usage demonstration, not an all-language model ranking. Translation and constraint
validity beyond the languages inspected remain important limitations.

\end{abstract}


\section{Introduction}
\label{sec:intro}

Instruction following evaluation has become central to developing practical AI systems
capable of understanding and executing natural language commands. The release of IFEval
\cite{zhou2023ifeval} established a standardised framework for measuring how well models
can follow verifiable instructions in English, providing a rigorous benchmark that goes
beyond traditional task completion metrics. However, IFEval's English-only scope reveals
a critical limitation: as large language models are increasingly deployed across diverse
linguistic contexts, evaluation benchmarks confined to English fail to capture
performance disparities that arise from multilingual instruction following.

Existing multilingual instruction-following datasets cover fewer languages:
M-IFEval covers French, Japanese, and Spanish \cite{dussolle-etal-2025-ifeval};
XIFBench covers six languages \cite{li2025xifbench}; and Marco-Bench-MIF covers
30 \cite{zeng-etal-2025-marco}. These resources differ in their construction and
evaluation procedures. Wider coverage creates new opportunities for analysis, but
translated constraints can also become invalid or unnatural.

We introduce \textbf{MultiIFEval}, extending the English IFEval prompts and their
verifiable constraints to 305 language configurations. The dataset is already public on Hugging Face. The reviewer-facing access link
must be made anonymous before submission; the non-anonymous release identifier
is recorded in the companion project notes, not this review manuscript.
It contains test-only language subsets, including languages with very different
amounts of available training data. Wikipedia text is supplied as translation
context, not as evidence that the resulting examples are culturally correct.

This paper contributes (1) a released multilingual resource and its construction
pipeline, and (2) a reproducible setup for a small demonstration with existing
instruction-following evaluation tooling. We deliberately separate breadth of
\emph{dataset coverage} from breadth of \emph{model evaluation}: an exhaustive
305-language leaderboard is beyond the scope of this submission.


\section{Related Work}
\label{sec:related-work}

\paragraph{IFEval}
% Original English benchmark
\cite{zhou2023ifeval} introduced IFEval, the first systematic benchmark for instruction-following evaluation in large language models. IFEval comprises approximately 500 prompts with verifiable constraints that can be checked programmatically, enabling objective binary evaluation of whether models successfully follow instructions. The benchmark covers four constraint types: content inclusion/exclusion, formatting requirements, length constraints, and punctuation rules. While IFEval established a rigorous framework for measurable instruction-following assessment, its scope is limited to English, leaving open the question of how models perform when following instructions in other languages.

\paragraph{M-IFEval}
% 3 languages, translated + language-specific
\cite{dussolle-etal-2025-ifeval} extended IFEval to three languages---French, Japanese, and Spanish---creating M-IFEval. This work evaluated both translated versions of the original English prompts and language-specific instructions designed natively in each target language. Their evaluation of 8 LLMs revealed notable cross-lingual performance gaps, with models generally performing worse on non-English instructions even when capabilities appeared comparable on English tasks. M-IFEval demonstrated that instruction-following quality does not transfer uniformly across languages, motivating broader multilingual evaluation.

\paragraph{XIFBench}
% 6 languages, constraint-based taxonomy
\cite{li2025xifbench} introduced XIFBench, expanding coverage to 6 languages with a fine-grained constraint taxonomy organised into five categories: Content, Style, Situation, Format, and Numerical constraints. XIFBench comprises 558 instructions designed to test specific constraint types, enabling diagnostic analysis of where models fail. Their taxonomy provides a structured framework for understanding instruction-following weaknesses, though coverage remains limited to a small set of predominantly high-resource languages.

\paragraph{Multi-IF}
% 8 languages, multi-turn
\cite{he2024multi} further expanded multilingual instruction following to 8 languages with Multi-IF, introducing the critical dimension of multi-turn interaction. Rather than evaluating single-turn instructions, Multi-IF assesses how well models follow sequences of constraints across multiple conversational turns. Their findings showed that models struggle significantly more with multi-turn constraints compared to single-turn instructions, with performance degrading as the number of constraints accumulates across turns. This work highlighted that real-world instruction following often requires maintaining constraints over extended interactions, not just single prompts.

\paragraph{Marco-Bench-MIF}
% 30 languages, localisation vs translation
\cite{zeng-etal-2025-marco} study translation and localisation with
Marco-Bench-MIF in 30 languages. MultiIFEval prioritises wider language coverage
rather than claiming comparable human validation across all languages.


\section{Dataset Generation}
\label{sec:dataset-generation}

For each English IFEval example, we prompt
Gemini-3-flash-preview~\cite{google-ai2026gemini-3-flash-preview} to translate the
prompt and localise language-dependent constraint values. Before generation,
English-specific constraints are adapted or removed where unsupported; the
resulting \texttt{instruction\_id\_list} is then held fixed. The zero-shot
prompt includes
a randomly selected Wikipedia article in the target language as contextual text;
this may help localisation, but does not guarantee factual or cultural accuracy.
Figure~\ref{fig:prompt} shows the generation prompt.

The generation schema represents \texttt{kwargs} as lists of name--value pairs
because the model's structured-output schema did not support dictionary-valued
fields. We convert these back to the evaluator's expected constraint metadata.
A valid JSON shape alone cannot establish that the translated prompt and its
machine-checkable constraints still mean the same thing.


\begin{figure*}[h]
\centering
\small
\begin{verbatim}
You are a professional translator from English to {target_language} (language code:
'{target_language_code}').

Here is an instruction-following example in English:

<example>
{"prompt": <prompt>, "instruction_id_list": <instruction_id_list>, "kwargs": <kwargs>}
</example>

You need to translate the example to {target_language}. This means the following:

1. The `prompt` should be translated to {target_language}. You should localise these to
   {target_language} and its country's culture.
2. The `instruction_id_list` should remain exactly the same.
3. The keys in the kwargs dictionaries should be completely unchanged, and the values
   should _only_ be changed if they contain English words or phrases, which would need
   to be changed to the equivalent in {target_language}.

Here is an example of some text written in {target_language}:

<{target_language_code}_example>
{wikipedia_article_in_target_language}
</{target_language_code}_example>

You should return the translated example in JSON format, with the following structure:

- `new_prompt` (str): The translated prompt.
- `new_instruction_id_list` (list[str]): The unchanged input instruction IDs.
- `new_kwargs` (list[list[dict]]): The translated keyword arguments.

Here the `new_kwargs` must be a list of the same length as the
`new_instruction_id_list`, containing lists of dicts, each with keys `name` and `value`,
where `name` is the name of the keyword argument and `value` is the value of the keyword
argument. Each list of kwarg dicts should be the keyword arguments for the corresponding
instruction ID in the `new_instruction_id_list`.
\end{verbatim}
\caption{Prompt used to generate the MultiIFEval dataset.}
\label{fig:prompt}
\end{figure*}


\subsection{Manual Danish Translation}
\label{sec:manual-danish-translation}

Two native Danish speakers also prepared a manual Danish translation of IFEval.
This provides a potential reference for a targeted translation-quality check,
subject to alignment of prompt IDs, availability of the manual data, and a
reported validation protocol. We do not treat a manual translation as an
error-free gold standard or assert unreported human-validation results.

\subsection{Machine Translated Dataset}
\label{sec:mt-dataset}

The public Hub release has 305 language configurations with a \texttt{test}
split and the columns \texttt{key}, \texttt{prompt},
\texttt{instruction\_id\_list}, and \texttt{kwargs}. Each config contains a
subset of the original prompts: generation or validation failures may remove
individual examples. The number of examples per language and the handling of
missing examples should be reported from a pinned release before submission.
The dataset's licence is CC BY-NC-SA 4.0. Automatic structure and ID checks do
not substitute for assessing semantic preservation or native-speaker quality.


\section{Focused Evaluation}
\label{sec:evaluation}

We plan a deliberately small demonstration of the released resource rather than a
full 305-language model survey. The initial comparison uses Danish and English
as a reference pair; German and French can be added if compute permits. We
will evaluate one or two small, instruction-tuned, non-reasoning decoder models
(e.g., Qwen2.5-1.5B-Instruct and Qwen2.5-3B-Instruct), with no few-shot examples.
The final submission must name exact model and dataset revisions, inference
settings, sample counts, and hardware. We do not use encoder-only models for
generative instruction following.

\paragraph{Reproducibility}
We will use EuroEval's existing instruction-following task with \emph{local}
\texttt{DatasetConfig} objects referencing the public language-specific Hub
configs; no changes to EuroEval are required. Its checker consumes the
\texttt{prompt}, \texttt{instruction\_id\_list}, and \texttt{kwargs} fields.
The local configuration and example invocation are included with this paper.
We will report EuroEval's instruction-accuracy metric, the number of evaluated
examples and applicable constraints, and any scorer incompatibilities or
language-dependent checking errors. Scores for these few languages cannot
support conclusions about the remaining language configurations.

\paragraph{Planned results (not yet measured)}
Insert a table of actual per-model, per-language scores with sample sizes,
followed by a short error analysis of constraint mismatches and translation
artefacts. Do not infer cross-lingual model rankings from unrun experiments.
If the manual Danish dataset and a documented assessment are available, add
an aligned, small qualitative comparison; otherwise omit human-validation
claims. The published English and Danish configurations are not themselves
manual-versus-machine parallel evaluations.

\section{Limitations and Ethical Considerations}
\label{sec:limitations}

The broad language coverage comes from automated translation, not equivalent
native-speaker validation in 305 languages. Wikipedia context may not match a
prompt's subject, and translated constraints (e.g., word counts or keywords)
may change difficulty or even become impossible. Availability of a language
configuration does not establish evaluation validity for that language.
The proposed demonstration covers only a few European languages and small
non-reasoning models; its findings must not be generalised to all language
communities or model families. The released dataset carries a
CC BY-NC-SA 4.0 licence; downstream users should review its terms.

\section{Conclusion}
\label{sec:conclusion}

MultiIFEval makes IFEval-style prompts and machine-checkable constraints
available in 305 language configurations. We document the public release and
an intentionally limited evaluation route. Broader model evaluation and
language-specific human quality assessment remain future work.

% Add authors, affiliations and acknowledgements only in the camera-ready version.
\bibliography{custom}

\appendix
\section{Evaluation Configuration}
\label{sec:appendix}
The companion local configuration demonstrates how to select a released
language subset and invoke EuroEval without registering a dataset upstream.
Replace this appendix with additional error examples or validation details
if they are available by submission.

\end{document}
